\chapter*{Glossaire}
\addcontentsline{toc}{chapter}{Glossaire}

\begin{description}[leftmargin=4cm, style=nextline, font=\color{BN}\bfseries]

  \item[Algorithme]
    Suite finie et non ambigu\"e d'instructions permettant de
    r\'esoudre un probl\`eme. Propri\'et\'es~: finitude, validit\'e,
    d\'eterminisme. Le mot vient du nom du math\'ematicien
    al-Khw\^arizm\^i ($\approx$820).

  \item[Affectation]
    Instruction $X\leftarrow\text{expr}$ qui stocke une valeur
    dans une variable, \'ecrasant l'ancienne. Ne pas confondre
    avec l'\'egalit\'e math\'ematique.

  \item[Dichotomie]
    M\'ethode de recherche ou d'approximation qui divise le
    domaine en deux \`a chaque \'etape.

  \item[Invariant de boucle]
    Propri\'et\'e logique vraie avant la boucle et conserv\'ee
    \`a chaque it\'eration~; permet de prouver la correction.

  \item[It\'eration]
    R\'ep\'etition d'un bloc d'instructions~: boucle \texttt{Pour}
    (nombre d'it\'erations connu) ou \texttt{Tant que}
    (arr\^et conditionnel).

  \item[Organigramme]
    Repr\'esentation graphique d'un algorithme.
    Conventions~: losange (test), trap\`eze (E/S),
    rectangle (traitement), fl\`eches (flux).

  \item[Pseudo-code]
    Langage informel, \`a mi-chemin entre le fran\c cais et un
    langage de programmation, utilis\'e pour d\'ecrire un algorithme
    ind\'ependamment de toute syntaxe particuli\`ere.

  \item[Tri]
    Algorithme qui r\'eordonne les valeurs d'un tableau, par exemple
    dans l'ordre croissant ou décroissant.

  \item[Variable]
    Emplacement m\'emoire nomm\'e pouvant contenir une valeur.
    En algorithmique, toute variable est d'un type
    (entier, r\'eel, bool\'een, cha\^ine\ldots).
    Structure alg\'ebrique munie de deux op\'erations ($\land$, $\lor$)
    et d'un compl\'ement ($\lnot$), satisfaisant des axiomes analogues
    \`a ceux des ensembles. Fond\`ee par George Boole (1854).

  \item[Assertion / Proposition]
    \'Enonc\'e ayant une valeur de v\'erit\'e d\'etermin\'ee~:
    vrai (\vrai) ou faux (\faux). Contraste avec le pr\'edicat
    (dont la valeur d\'epend d'une variable).

  \item[Condition n\'ecessaire]
    $P$ est \emph{condition n\'ecessaire} pour $Q$ si
    $Q\Rightarrow P$. Autrement dit~: sans $P$, pas de $Q$.

  \item[Condition suffisante]
    $P$ est \emph{condition suffisante} pour $Q$ si $P\Rightarrow Q$.
    Autrement dit~: $P$ \og suffit\fg\ pour garantir $Q$.

  \item[Connecteur logique]
    Symbole op\'erant sur une ou plusieurs propositions~:
    $\lnot$ (n\'egation), $\land$ (conjonction), $\lor$ (disjonction),
    $\Rightarrow$ (implication), $\Leftrightarrow$ (\'equivalence).

  \item[Contradiction]
    Proposition \emph{toujours fausse}, quelle que soit la valeur
    de ses variables (ex.~: $P\land\lnot P$).

  \item[Contrapos\'ee]
    La contrapos\'ee de $P\Rightarrow Q$ est $\lnot Q\Rightarrow\lnot P$.
    Une implication et sa contrapos\'ee sont toujours \'equivalentes.

  \item[Contre-exemple]
    Un \'el\'ement $x_0$ tel que $P(x_0)$ est fausse, utilis\'e pour
    r\'efuter $\forall x,\,P(x)$.

  \item[D\'emonstration par l'absurde]
    Raisonnement dans lequel on suppose $\lnot P$ et l'on aboutit \`a
    une contradiction, concluant que $P$ est vraie.

  \item[Discriminant]
    Pour $ax^2+bx+c=0$~: $\Delta=b^2-4ac$.
    Son signe d\'etermine le nombre de racines r\'eelles.

  \item[\'Equivalence logique]
    $P\Leftrightarrow Q$~: vraie si $P$ et $Q$ ont la m\^eme valeur
    de v\'erit\'e. \'Equivalent \`a $(P\Rightarrow Q)\land(Q\Rightarrow P)$.

  \item[Implication]
    $P\Rightarrow Q$~: fausse uniquement quand $P=\vrai$ et $Q=\faux$.
    Se lit \og si $P$ alors $Q$\fg.

  \item[Irrationnel]
    Nombre r\'eel qui ne peut pas s'\'ecrire $p/q$ avec $p\in\Z$,
    $q\in\N^*$. Exemples~: $\sqrt{2}$, $\sqrt{3}$, $\pi$, $e$.

  \item[Lois de De Morgan]
    $\lnot(P\land Q)\equiv\lnot P\lor\lnot Q$ ;
    $\lnot(P\lor Q)\equiv\lnot P\land\lnot Q$.
    Permettent de distribuer la n\'egation.

  \item[N\'egation]
    $\lnot P$~: vraie si $P$ est fausse, fausse si $P$ est vraie.

  \item[Pr\'edicat]
    \'Enonc\'e $P(x)$ dont la valeur de v\'erit\'e d\'epend de $x$.
    Devient une proposition pour une valeur particuli\`ere de $x$.

  \item[Quantificateur existentiel $\exists$]
    $\exists x\in E,\,P(x)$~: il existe au moins un $x$ dans $E$
    pour lequel $P(x)$ est vraie.

  \item[Quantificateur universel $\forall$]
    $\forall x\in E,\,P(x)$~: pour tout $x$ dans $E$,
    $P(x)$ est vraie.

  \item[R\'eciproque]
    La r\'eciproque de $P\Rightarrow Q$ est $Q\Rightarrow P$.
    Elle n'est \textbf{pas} \'equivalente \`a l'implication initiale
    en g\'en\'eral.

  \item[Table de v\'erit\'e]
    Tableau listant la valeur d'une formule logique pour toutes
    les combinaisons de valeurs de ses variables.

  \item[Tautologie]
    Proposition \emph{toujours vraie}, quelle que soit la valeur
    de ses variables (ex.~: $P\lor\lnot P$).

  \item[Vecteur]
    Objet d\'efini par une direction, un sens et une norme.
    Dans un rep\`ere~: $\vec u=(x_1,y_1)$.

\end{description}

%% ── NOTATIONS ───────────────────────────────────────────────────
